\begin{abstract}
Parthenium hysterophorus is an aggressively expanding invasive weed in South Asia, yet its local occurrence patterns, climatic suitability, and overlap with the biocontrol beetle \textit{Zygogramma bicolorata} remain poorly resolved. This study investigated habitat, road, and recent climate factors associated with \textit{Parthenium} presence in Pakistan; projected changes in climatic suitability across India under future climate scenarios; and the concordance between a Pakistan-calibrated generalised additive model (GAM) and an India-calibrated MaxEnt model in Pakistan. In Pakistan, roadside presence--absence surveys analysed via a binomial GAM revealed that \textit{Parthenium} occurrence was highest in open and disturbed habitats and varied by road class and recent climate, with residual spatial structure persisting. In India, a presence-only MaxEnt model projected extensive current climatic suitability for \textit{Parthenium}, but predicted substantial restriction and fragmentation by 2030 and 2050. Notably, overlap with a temperature-only annual thermal-permissiveness proxy for \textit{Z. bicolorata} declined from 66.0\% under current conditions to 42.0\% in 2030 and 20.6\% in 2050, largely due to contraction in \textit{Parthenium} suitability. Cross-model comparison across Pakistan indicated only weak concordance, highlighting model-dependent uncertainty. These findings underscore the need for cautious, spatially targeted monitoring and management in the face of climate change.
\end{abstract}